\documentclass[11pt]{article}

\usepackage[T1]{fontenc}
\usepackage[utf8]{inputenc}
\usepackage[spanish]{babel}
\usepackage{amsmath,amssymb}
\usepackage{booktabs}
\usepackage{geometry}
\usepackage{hyperref}
\usepackage{enumitem}
\usepackage{listings}
\geometry{margin=2.5cm}
\lstset{basicstyle=\ttfamily\small, breaklines=true}

\title{Auditor\'ia T\'ecnica del Pipeline GNN (Tab a Tab)\\
\large Documentaci\'on detallada, GraphSMOTE e ImGAGN}
\author{}
\date{\today}

\begin{document}
\maketitle
\tableofcontents

\section{Alcance y metodolog\'ia}
Este informe documenta y audita el pipeline GNN implementado en el repositorio, con foco en:
\begin{itemize}[leftmargin=*]
  \item El flujo \emph{tab a tab} de la aplicaci\'on Streamlit (\texttt{src/graph\_builder\_app.py}).
  \item La construcci\'on del grafo, el modelo GNN, el entrenamiento y la evaluaci\'on (\texttt{src/graph.py}, \texttt{src/gat\_model.py}, \texttt{src/gnn\_main.py}, \texttt{src/train\_pretrain.py}).
  \item Los mecanismos de balanceo: GraphSMOTE e ImGAGN (\texttt{src/graphsmote.py}, \texttt{src/imgagn.py}).
  \item La b\'usqueda de hiperpar\'ametros (Optuna) y el control de m\'etricas.
\end{itemize}
La auditor\'ia incluye verificaci\'on de mejores pr\'acticas, riesgos de fuga de informaci\'on, uso de memoria/c\'alculo, y consistencia con las opciones de configuraci\'on (\texttt{src/config.py}).

\section{Mapa general del pipeline}
\textbf{Resumen conceptual:}
\begin{enumerate}[leftmargin=*]
  \item \textbf{Datos base}: p\'orticos, flujos y accidentes (\texttt{src/utils.py}).
  \item \textbf{Feature engineering}: snapshot/flow \textrm{(\texttt{src/snapshot\_pipeline.py}, \texttt{src/features.py})}.
  \item \textbf{Feature selection}: selecci\'on manual/guiada (UI) y persistencia.
  \item \textbf{Grafo}: nodos \texttt{pm}, aristas \texttt{temporal/spatial/st\_fwd} y atributos (\texttt{src/graph.py}).
  \item \textbf{Entrenamiento}: HeteroGAT + TemporalAggregator (opcional) (\texttt{src/gat\_model.py}, \texttt{src/temporal\_head.py}).
  \item \textbf{Balanceo}: GraphSMOTE (embedding-space + z\textrm{\textrightarrow}x + edge-gen) y/o ImGAGN (adversarial) (\texttt{src/graphsmote.py}, \texttt{src/imgagn.py}).
  \item \textbf{Evaluaci\'on}: m\'etricas (F1, AUC, AUPRC, MCC), umbrales y calibraci\'on (\texttt{src/gnn\_main.py}).
  \item \textbf{Artefactos}: modelos, hparams JSON, grafos aumentados, historial (\texttt{Resultados/}, \texttt{gnn\_history.jsonl}).
\end{enumerate}

\section{Pipeline \emph{tab a tab} (Streamlit Graph Builder)}
El flujo UI se implementa en \texttt{src/graph\_builder\_app.py} y organiza el trabajo en tabs.

\subsection{Eventos}
\textbf{Funci\'on UI}: \texttt{\_render\_events\_tab}.\\
\textbf{Objetivo}: cargar y validar accidentes, normalizar columnas, limpiar y preparar el dataset de eventos.
\begin{itemize}[leftmargin=*]
  \item Carga de accidentes desde CSV/archivos seleccionados, persistencia en \texttt{st.session\_state}.
  \item Uso de utilidades en \texttt{src/utils.py} para normalizar columnas y fechas.
  \item Soporte para generar reportes de \emph{Puntos Negros} (posteriormente usados en \emph{Graph}).
\end{itemize}
\textbf{Implicaci\'on}: este tab define el universo de accidentes, que afectar\'a etiquetas futuras y cortes temporales.

\subsection{Feature Engineering}
\textbf{Funci\'on UI}: \texttt{\_render\_feature\_engineering}.\\
\textbf{Objetivo}: generar o cargar features por nodo \texttt{pm}.
\begin{itemize}[leftmargin=*]
  \item Modos principales: \emph{Snapshot Features} y \emph{Flow 5 min} (ver \texttt{src/snapshot\_pipeline.py}).
  \item Generaci\'on de features con \texttt{compute\_pm\_features} (\texttt{src/features.py}) y/o \texttt{SnapshotFeatureBuilder}.
  \item Salida persistida en DuckDB o CSV; sincronizaci\'on de rango temporal y tramo.
\end{itemize}
\textbf{Implicaci\'on}: afecta dimensionalidad de nodos y distribuci\'on de atributos. La resoluci\'on temporal se gobierna con \texttt{DT\_MINUTES} y par\'ametros en \texttt{src/config.py}.

\subsection{Feature Explorer}
\textbf{Funci\'on UI}: \texttt{render\_feature\_explorer}.\\
\textbf{Objetivo}: inspecci\'on visual, estad\'isticas y outliers de variables de entrada.
\textbf{Implicaci\'on}: diagn\'ostico de calidad de features antes de construir el grafo.

\subsection{Feature Selection}
\textbf{Funci\'on UI}: \texttt{\_render\_feature\_selection\_tab}.\\
\textbf{Objetivo}: seleccionar subconjuntos de variables para entrenamiento, guardar listas y metadatos.
\textbf{Implicaci\'on}: cambia el espacio de entrada del GNN y la dimensionalidad de \texttt{pm.x}.

\subsection{Graph}
\textbf{Funciones UI}: \texttt{\_render\_create\_graph}, \texttt{\_render\_load\_graph}, \texttt{\_render\_in\_memory\_graph}.\\
\textbf{Objetivo}: construir o cargar el grafo heterog\'eneo de p\'orticos.
\begin{itemize}[leftmargin=*]
  \item Selecci\'on de tramo (manual o por puntos negros).
  \item Construcci\'on de nodos \texttt{pm} (p\'ortico-minuto) con features normalizados.
  \item Generaci\'on de aristas temporales/espaciales y atributos de arista (ver Secci\'on \ref{sec:grafo}).
  \item Creaci\'on de \texttt{train/val/test\_mask} con corte temporal (\texttt{\_compute\_time\_cutoffs}).
  \item Guardado en \texttt{Resultados/highway\_graph\_*.pt} y en memoria (\texttt{st.session\_state}).
\end{itemize}

\subsection{Visual Graph}
\textbf{Funci\'on UI}: \texttt{render\_visual\_graph\_tab}.\\
\textbf{Objetivo}: visualizar topolog\'ia, grados, y subgrafos para QA.

\subsection{Network}
\textbf{Funci\'on UI}: \texttt{\_render\_network\_tab}.\\
\textbf{Objetivo}: configurar arquitectura y par\'ametros del modelo antes de entrenar.
\textbf{Implicaci\'on}: fija opciones como \texttt{hidden\_channels}, \texttt{num\_heads}, \texttt{dropout}, \texttt{num\_layers}, agregaciones \texttt{aggr1/aggr2} y checkpointing.

\subsection{Optuna}
\textbf{Funci\'on UI}: \texttt{\_render\_optuna\_tab}.\\
\textbf{Objetivo}: b\'usqueda de hiperpar\'ametros con Optuna (ver Secci\'on \ref{sec:optuna}).
\textbf{Implicaci\'on}: almacena \texttt{optuna\_hyperparams\_*.csv} y gobierna la selecci\'on autom\'atica en entrenamiento.

\subsection{Balance}
\textbf{Funci\'on UI}: \texttt{\_render\_balance\_tab}.\\
\textbf{Objetivo}: generar grafos balanceados con GraphSMOTE o ImGAGN.
\begin{itemize}[leftmargin=*]
  \item GraphSMOTE: requiere modelo entrenado para embeddings y decodificadores \texttt{z\textrightarrow x}.
  \item ImGAGN: entrena generador y discriminador adversarial y crea nodos sint\'eticos.
\end{itemize}
\textbf{Implicaci\'on}: genera grafos con \texttt{is\_synthetic}, actualiza \texttt{train\_mask}, y puede cambiar el balance real.

\subsection{Training}
\textbf{Funci\'on UI}: \texttt{\_render\_training\_tab}.\\
\textbf{Objetivo}: entrenar el GNN con opci\'on de grafo balanceado, early stopping y evaluaci\'on autom\'atica.
\textbf{Backend}: \texttt{src/gnn\_main.py::run\_gat\_training} y \texttt{src/train\_pretrain.py::train\_minibatch}.

\subsection{Evaluaci\'on}
\textbf{Funci\'on UI}: \texttt{\_render\_evaluation\_tab}.\\
\textbf{Objetivo}: cargar modelos, evaluar m\'etricas, aplicar umbral y calibraci\'on opcional.
\textbf{Backend}: \texttt{src/gnn\_main.py::test} y utilidades de exportaci\'on.

\subsection{History}
\textbf{Funci\'on UI}: \texttt{\_render\_history\_tab}.\\
\textbf{Objetivo}: registro de experimentos y trazabilidad (\texttt{Resultados/gnn\_history.jsonl}).

\subsection{Experiments}
\textbf{Funci\'on UI}: \texttt{\_render\_gnn\_experiments\_tab}.\\
\textbf{Objetivo}: ejecutar lotes experimentales, exportar resultados, e importar ZIPs.

\section{Modelo de datos del grafo}
\label{sec:grafo}
\textbf{Nodos}: \texttt{pm} (p\'ortico-minuto).\\
\textbf{Targets}: \texttt{pm.y} (etiqueta binaria), \texttt{pm.is\_accident\_pm} (accidente directo).\\
\textbf{M\'ascaras}: \texttt{train\_mask/val\_mask/test\_mask/sequence\_mask}.

\subsection{Aristas y atributos}
\textbf{Aristas principales} (\texttt{src/graph.py}):
\begin{itemize}[leftmargin=*]
  \item \texttt{temporal}: $P_{i,t} \to P_{i,t+1}$ (misma ubicaci\'on, siguiente tiempo).
  \item \texttt{spatial}: $P_{i,t} \to P_{i+1,t}$ (siguiente p\'ortico, mismo tiempo).
  \item \texttt{st\_fwd} (opcional): $P_{i,t} \to P_{i+1,t+1}$.
  \item \texttt{spatial\_back} (opcional): inversa de \texttt{spatial}.
\end{itemize}
\textbf{Atributos de arista}: diferencia de features entre destino y origen ($\Delta x$) y extras
para \texttt{spatial} (gradientes, shock, distancia). En \texttt{temporal} y \texttt{st\_fwd}
se usan \emph{extras} en cero para alinear dimensiones.

\subsection{Etiquetado temporal}
\textbf{Cortes temporales}: definidos por accidentes ordenados (\texttt{\_compute\_time\_cutoffs}).\\
\textbf{Mascaras}: aplicadas por timestamps en minutos.\\
\textbf{SequenceIndex}: secuencias temporales con guard band y horizonte futuro
(\texttt{src/snapshot\_sequences.py}).

\section{Arquitectura del modelo GNN}
\textbf{Modelo base}: \texttt{src/gat\_model.py::HeteroGAT}.
\begin{itemize}[leftmargin=*]
  \item \textbf{Capas}: \texttt{num\_layers} bloques \texttt{HeteroConv} con \texttt{GATConvSaveAlpha}.
  \item \textbf{Agregaci\'on}: \texttt{aggr1} en la primera capa, \texttt{aggr2} en capas siguientes.
  \item \textbf{Atenciones}: capturadas cuando \texttt{XAI=1}, usadas para regularizaci\'on $L2$.
  \item \textbf{Edge attributes}: se validan por batch; si faltan o desalinean, se rellenan con ceros.
  \item \textbf{Checkpointing}: \texttt{use\_checkpointing} reduce memoria con recomputaci\'on.
\end{itemize}

\subsection{Temporal head (opcional)}
\textbf{Clase}: \texttt{src/temporal\_head.py::TemporalAggregator}.\\
\textbf{Idea}: GRU sobre secuencias de embeddings recientes; mantiene cache de embeddings para nodos no presentes en el batch, evitando fuga de gradientes.

\subsection{Registro de variantes GNN (encoder + temporal + aristas)}
\label{sec:gnn_variantes}
El pipeline ahora incorpora un \textbf{registro de variantes} controlado por
\texttt{src/config.py} (\texttt{GNN\_VARIANT}, \texttt{GNN\_TEMPORAL\_CACHE\_STRATEGY})
y materializado en \texttt{src/gnn\_main.py} mediante las funciones
\texttt{\_build\_gnn\_model}, \texttt{\_build\_temporal\_head\_for\_variant} y
\texttt{\_prime\_temporal\_cache\_if\_needed}. Adem\'as, la UI en
\texttt{src/graph\_builder\_app.py} expone una gu\'ia resumida en los tabs
\emph{Network}, \emph{Training}, \emph{Optuna} y \emph{Evaluaci\'on}. La idea de dise\~no es separar:
\begin{itemize}[leftmargin=*]
  \item \textbf{Encoder espacial}: GAT heterog\'eneo base, GAT con codificador de aristas, o variante edge-aware.
  \item \textbf{Head temporal}: snapshot (sin secuencia), GRU, Transformer temporal o attention pooling.
  \item \textbf{Protocolo de cache}: \texttt{incremental} por batch o \texttt{global\_epoch} (forward global por \'epoca para secuencias consistentes).
\end{itemize}

\subsubsection{Idea general de comparabilidad}
Todas las variantes comparten el mismo \texttt{SequenceIndex} (construido en
\texttt{src/snapshot\_sequences.py} a partir de \texttt{sequence\_rows}/\texttt{target\_rows}),
el mismo grafo y los mismos splits temporales. Esto permite comparar arquitecturas sin cambiar
la definici\'on de etiqueta ni el horizonte de predicci\'on.

\paragraph{Etiquetado ``5 minutos antes''}
La l\'ogica de etiquetado usa:
\[
\text{intervalo positivo} = (t+\gamma,\; t+\gamma+H]
\]
donde \(\gamma\) es \texttt{GUARD\_BAND\_MINUTES} y \(H\) es \texttt{HORIZON\_MINUTES}.
Para una configuraci\'on estricta de ``5 minutos antes'', se recomienda:
\begin{itemize}[leftmargin=*]
  \item \texttt{GUARD\_BAND\_MINUTES = 5}
  \item \texttt{HORIZON\_MINUTES = 5} (o 10--20 para riesgo en ventana)
\end{itemize}
en \texttt{src/config.py}.

\subsubsection{Variantes implementadas}
\textbf{Implementaci\'on}: \texttt{src/gat\_model.py}, \texttt{src/temporal\_head.py},
\texttt{src/temporal\_heads.py}, \texttt{src/gnn\_main.py}.\\
\textbf{Selecci\'on}: mediante el string \texttt{GNN\_VARIANT} (normalizado en
\texttt{src/gnn\_main.py::\_normalize\_gnn\_variant}).

\begin{enumerate}[leftmargin=*]
  \item \textbf{\texttt{gat\_snapshot}} (baseline): \texttt{HeteroGAT} sin head temporal. Clasifica con el embedding del snapshot actual.
  \item \textbf{\texttt{gat\_gru}}: \texttt{HeteroGAT} + \texttt{TemporalGRUHead} (\texttt{src/temporal\_heads.py}).
  \item \textbf{\texttt{gat\_transformer}}: \texttt{HeteroGAT} + \texttt{TemporalTransformerHead} (self-attention con positional embedding aprendible).
  \item \textbf{\texttt{gat\_attnpool}}: \texttt{HeteroGAT} + \texttt{TemporalAttnPoolHead} (ablation ligera y barata computacionalmente).
  \item \textbf{\texttt{gat\_edge\_mlp}}: \texttt{HeteroGATWithEdgeEncoder} (MLP por tipo de arista sobre \texttt{edge\_attr}) sin head temporal.
  \item \textbf{\texttt{gat\_edge\_mlp\_gru}} y \textbf{\texttt{gat\_edge\_mlp\_transformer}}: combinan codificador de aristas + head temporal.
  \item \textbf{\texttt{gnn\_edge\_aware}}: \texttt{HeteroEdgeAware} (usa \texttt{TransformerConv}) sin head temporal.
  \item \textbf{\texttt{gnn\_edge\_aware\_gru}} y \textbf{\texttt{gnn\_edge\_aware\_transformer}}: variante edge-aware con modelado temporal.
\end{enumerate}

\subsubsection{Codificaci\'on de aristas (\texttt{edge\_attr})}
La variante \texttt{HeteroGATWithEdgeEncoder} agrega un MLP de aristas
(\texttt{EdgeAttrEncoder}) antes de la atenci\'on. Motivaci\'on:
\begin{itemize}[leftmargin=*]
  \item los deltas \(\Delta x\) pueden tener escalas heterog\'eneas;
  \item un mapeo aprendido puede inducir una geometr\'ia m\'as estable para la atenci\'on;
  \item permite estudiar de forma separada el efecto de ``mejor \texttt{edge\_attr}'' vs. ``mejor head temporal''.
\end{itemize}

\subsubsection{Heads temporales y ensamblado de secuencias}
Los heads temporales reutilizan la infraestructura de \texttt{TemporalAggregator}:
\begin{itemize}[leftmargin=*]
  \item \texttt{update\_cache(...)} para cache global o parcial de embeddings.
  \item \texttt{\_build\_sequence\_batch(...)} para reconstruir tensores \([B, L, D]\) usando \texttt{sequence\_rows}.
\end{itemize}
Esto permite cambiar el operador temporal (GRU/Transformer/Attention Pooling) sin cambiar el
contrato del entrenamiento en \texttt{src/train\_pretrain.py} ni la firma de \texttt{forward}.

\subsubsection{Estrategias de cache temporal (punto cr\'itico)}
\begin{itemize}[leftmargin=*]
  \item \textbf{\texttt{incremental}}: cache por batch (hist\'orico, menor costo, m\'as dependiente del sampler).
  \item \textbf{\texttt{global\_epoch}}: priming por \'epoca con forward global (\texttt{compute\_epoch\_embeddings}), luego training/eval usan cache consistente.
\end{itemize}
\textbf{Recomendaci\'on para art\'iculo/ciencia}: usar \texttt{global\_epoch} para que las secuencias
temporales tengan la misma referencia de embeddings dentro de una \'epoca.

\subsubsection{Matriz experimental recomendada}
\begin{enumerate}[leftmargin=*]
  \item \texttt{gat\_snapshot} (baseline espacial)
  \item \texttt{gat\_gru}
  \item \texttt{gat\_transformer}
  \item \texttt{gat\_edge\_mlp\_gru}
  \item \texttt{gnn\_edge\_aware\_gru} (o \texttt{gnn\_edge\_aware\_transformer})
\end{enumerate}
Mantener semilla, split temporal, \texttt{SequenceIndex} y protocolo de evaluaci\'on constantes.

\subsubsection{M\'etricas para eventos raros}
Para comparaci\'on cient\'ifica y operativa, priorizar:
\begin{itemize}[leftmargin=*]
  \item \textbf{PR-AUC / AUPRC}: m\'etrica principal en desbalance severo.
  \item \textbf{ROC-AUC}: referencia secundaria.
  \item \textbf{Recall@FAR/FPR objetivo}: util para restricciones operacionales.
  \item \textbf{MCC}: robusta cuando hay fuerte asimetr\'ia de clases.
\end{itemize}

\subsubsection{Trazabilidad de variantes}
El entrenamiento y HPO guardan \texttt{gnn\_variant} y \texttt{variant\_tag}
en sidecars de hiperpar\'ametros, checkpoints y logs estructurados
(\texttt{Resultados/gat\_model\_BEST*.pt}, \texttt{\_hparams.json}, Optuna CSVs).
Esto evita mezclar checkpoints de arquitecturas incompatibles.

\section{Entrenamiento y optimizaci\'on}
\textbf{Punto de entrada}: \texttt{src/gnn\_main.py::run\_gat\_training}.

\subsection{Flujo del entrenamiento}
\begin{enumerate}[leftmargin=*]
  \item \textbf{Dispositivo}: \texttt{get\_auto\_device()} (MPS $>$ CUDA $>$ CPU).
  \item \textbf{Carga de datos}: \texttt{loaded\_obj['data']}.
  \item \textbf{ImGAGN autom\'atico} (si \texttt{AUTO\_IMGAGN\_PRETRAIN=1}).
  \item \textbf{GraphSMOTE}: elecci\'on por usuario/flag y detecci\'on de sint\'eticos existentes.
  \item \textbf{HPO}: selecci\'on de \texttt{optuna\_hyperparams\_*.csv} o b\'usqueda nueva.
  \item \textbf{Modelo}: f\'abrica \texttt{\_build\_gnn\_model} selecciona encoder/temporal head seg\'un \texttt{gnn\_variant} y adjunta el head temporal cuando aplica.
  \item \textbf{Decodificadores z\textrightarrow x}: entrenados si GraphSMOTE activo (\texttt{train\_z2x\_decoders}).
  \item \textbf{Edge generator}: \texttt{RelEdgeGen} para reconstrucci\'on y aristas sint\'eticas.
  \item \textbf{P\'erdida}: \texttt{CrossEntropy} o \texttt{FocalLoss} (con/ sin pesos).
  \item \textbf{Scheduler}: \texttt{OneCycleLR}.
  \item \textbf{Loop}: entrenamiento por mini-batches, validaci\'on, early stopping, guardado de mejor modelo.
\end{enumerate}

\subsection{Train-minibatch y regularizaci\'on}
\textbf{Funci\'on}: \texttt{src/train\_pretrain.py::train\_minibatch}.
\begin{itemize}[leftmargin=*]
  \item Clasificaci\'on sobre \texttt{pm} (primeros \texttt{batch\_size} del \texttt{NeighborLoader}).
  \item P\'erdida total:
  \[
  \mathcal{L} = \mathcal{L}_{cls} + \lambda_{edge}\,\mathcal{L}_{edge} + \lambda_{att}\,\mathcal{L}_{att}.
  \]
  \item \textbf{\(\mathcal{L}_{edge}\)}: reconstrucci\'on de aristas con muestreo negativo (si \texttt{edge\_gen}).
  \item \textbf{\(\mathcal{L}_{att}\)}: $L2$ sobre atenciones cuando \texttt{XAI=1}.
  \item Gradiente acumulado (\texttt{accumulation\_steps=2}) y \texttt{grad\_clip}.
  \item Soporte AMP (solo CUDA).
\end{itemize}

\section{GraphSMOTE: s\'intesis de nodos y aristas}
\label{sec:graphsmote}
\textbf{M\'odulo}: \texttt{src/graphsmote.py}.\\
\textbf{Principio}: generar nodos sint\'eticos en el espacio embedding y decodificarlos a features reales.

\subsection{Paso 1: embeddings}
\begin{itemize}[leftmargin=*]
  \item \texttt{get\_embeddings\_minibatch} (para balance offline en UI) o \texttt{compute\_train\_embeddings} (en entrenamiento, sin fuga).
  \item Normalizaci\'on $\ell_2$ por nodo.
\end{itemize}

\subsection{Paso 2: SMOTE en embeddings}
Se aplica SMOTE en el conjunto minoritario:
\[
\mathbf{z}_{syn} = (1-\alpha)\,\mathbf{z}_i + \alpha\,\mathbf{z}_j,\quad \alpha \sim U(0,1).
\]
\begin{itemize}[leftmargin=*]
  \item k-NN sobre embeddings minoritarios (CPU cache o GPU si cabe).
  \item Par\'ametros clave: \texttt{GRAPHSMOTE\_K}, \texttt{smote\_k}, \texttt{random\_state}.
\end{itemize}

\subsection{Paso 3: decodificadores z\textrightarrow x}
\textbf{Clase}: \texttt{Z2XDecoders}.\\
\textbf{Entrenamiento}: \texttt{train\_z2x\_decoders} con SmoothL1/MSE, early stopping y validaci\'on.
\begin{itemize}[leftmargin=*]
  \item Un MLP por tipo de nodo: $\mathbf{x}\_{syn} = f_{ntype}(\mathbf{z}\_{syn})$.
  \item Se respetan \texttt{train/val\_mask} y se evita usar nodos fuera del split cuando est\'a disponible.
\end{itemize}

\subsection{Paso 4: inserci\'on de nodos}
\begin{itemize}[leftmargin=*]
  \item Concatenar \texttt{x} y \texttt{y} sint\'eticos a \texttt{pm.x} y \texttt{pm.y}.
  \item Actualizar \texttt{train/val/test\_mask} (por defecto nuevos nodos van a train).
  \item Crear/actualizar \texttt{is\_synthetic} y \texttt{is\_accident\_pm}.
\end{itemize}

\subsection{Paso 5: generaci\'on de aristas sint\'eticas}
\textbf{Generador}: \texttt{RelEdgeGen} (par\'ametro por relaci\'on).\\
\textbf{Candidatos reales}:
\begin{itemize}[leftmargin=*]
  \item \texttt{GRAPHSMOTE\_CONECT=1}: top-K en \texttt{train\_mask}.
  \item \texttt{GRAPHSMOTE\_CONECT=0}: vecinos a 1 salto de accidentes reales.
\end{itemize}
\textbf{Direcci\'on}:
\begin{itemize}[leftmargin=*]
  \item \texttt{BIDIRECCTION=1}: real $\leftrightarrow$ sint\'etico.
  \item \texttt{BIDIRECCTION=0}: solo real $\to$ sint\'etico.
\end{itemize}
\textbf{Atributos de arista}: se computan como $\Delta x$ entre nodos destino y origen (con \texttt{delta\_feature\_idx} si existe) y se alinean a la dimensionalidad existente.

\subsection{Modos de operaci\'on}
\begin{itemize}[leftmargin=*]
  \item \textbf{Offline}: \texttt{augment\_graph\_offline\_once} aumenta una vez y usa el grafo para entrenar.
  \item \textbf{Online}: \texttt{refresh\_synthetics\_online} regenera nodos cada \texttt{SMOTE\_EVERY\_N\_EPOCHS}.
  \item \textbf{None}: sin aumentaci\'on.
\end{itemize}

\section{ImGAGN: generaci\'on adversarial}
\label{sec:imgagn}
\textbf{M\'odulo}: \texttt{src/imgagn.py}.\\
\textbf{Idea central}: generar nodos sint\'eticos como combinaciones convexas de nodos minoritarios \emph{de entrenamiento} y entrenar un discriminador adversarial.

\subsection{Paso 1: homogenizaci\'on}
\texttt{to\_homogeneous\_if\_needed} transforma \texttt{HeteroData} a grafo homog\'eneo (target \texttt{pm}), preservando \texttt{edge\_attr} y \texttt{delta\_feature\_idx} si existen.

\subsection{Paso 2: definir cantidad de sint\'eticos}
\begin{align*}
\lambda_1 &= \frac{n_{min} + n_g}{n_{maj}} \\
\Rightarrow n_g &= \max\bigl(0,\, \lambda_1 n_{maj} - n_{min}\bigr)
\end{align*}
Se limita con \texttt{max\_new\_nodes} para evitar OOM.

\subsection{Paso 3: generador}
\textbf{Clase}: \texttt{ImGAGNGenerator}.\\
\begin{itemize}[leftmargin=*]
  \item Entrada: ruido $\mathbf{z} \sim \mathcal{N}(0, I)$.
  \item Salida: logits por nodo minoritario de entrenamiento.
  \item Pesos $T = \mathrm{softmax}(G(z))$ y features sint\'eticos $X_g = T X_{min}$.
  \item Aristas: top-K conexiones a nodos minoritarios reales seg\'un $T$.
\end{itemize}

\subsection{Paso 4: discriminador}
\textbf{Clase}: \texttt{ImGAGNDiscriminator} (GCN + 2 cabezas).
\begin{itemize}[leftmargin=*]
  \item \textbf{Head real/fake} (BCE).
  \item \textbf{Head minority/majority} (BCE).
  \item T\'ermino de separaci\'on \texttt{\_mm\_separation} para alejar medias minoritaria y mayoritaria.
\end{itemize}

\subsection{Paso 5: entrenamiento adversarial}
Por cada \texttt{epoch}:
\begin{enumerate}[leftmargin=*]
  \item Muestrear ruido y construir sint\'eticos (sin gradientes).
  \item Construir grafo aumentado y entrenar D por \texttt{d\_steps} minibatches.
  \item Recalcular $X_g$ con gradientes y optimizar G para:
  \[
  \mathcal{L}_G = \mathcal{L}_{rf} + \mathcal{L}_{mi} + \mathcal{L}_{di},
  \]
  donde $\mathcal{L}_{rf}$ empuja a \emph{real}, $\mathcal{L}_{mi}$ empuja a \emph{minority}, y $\mathcal{L}_{di}$ acerca embeddings a centroides minoritarios.
\end{enumerate}
\textbf{Escalabilidad}: usa \texttt{NeighborLoader} si est\'a disponible; si no, full-batch.

\subsection{Paso 6: grafo final}
Se reconstruye el grafo aumentado con el generador final y se devuelven:
\texttt{x\_aug}, \texttt{edge\_index\_aug}, \texttt{edge\_attr\_aug}, y la porci\'on de nodos generados.

\section{Optuna (b\'usqueda de hiperpar\'ametros)}
\label{sec:optuna}
\textbf{Funci\'on}: \texttt{src/gnn\_main.py::objective}.
\begin{itemize}[leftmargin=*]
  \item Objetivo: maximizar F0.5 en validaci\'on (prioriza precisi\'on).
  \item Espacio: \texttt{hidden\_channels}, \texttt{num\_heads}, \texttt{dropout}, \texttt{num\_layers}, \texttt{aggr1/aggr2}.
  \item Optimizadores: Adam/AdamW/RAdam (configurable).
  \item Con GraphSMOTE: busca \texttt{smote\_k}, \texttt{target\_pos\_ratio}, \texttt{smote\_every\_n\_epochs}, \texttt{lambda\_edge}.
  \item Sin GraphSMOTE: \texttt{FocalLoss} con \texttt{focal\_alpha}, \texttt{focal\_gamma}.
\end{itemize}
\textbf{Muestreo de seeds (train/val)}:
\begin{itemize}[leftmargin=*]
  \item \emph{Resamplear seeds por epoch}: cambia el subset en cada epoch (mide robustez).
  \item \emph{Fijar subset por trial}: mantiene el mismo subset durante el trial (reduce varianza).
  \item \emph{Cargar todo el grafo en memoria (sin muestreo)}: usa todos los nodos \texttt{train/val} (equivale a \texttt{sample\_train\_nodes=0} y \texttt{sample\_val\_nodes=0}); mayor uso de memoria.
\end{itemize}

\section{Evaluaci\'on, calibraci\'on y artefactos}
\begin{itemize}[leftmargin=*]
  \item \textbf{Evaluaci\'on}: \texttt{test} calcula F1, AUC, AUPRC, MCC, matriz de confusi\'on.
  \item \textbf{Umbral}: \texttt{pick\_tau\_fbeta} busca umbral \(\tau\) en validaci\'on.
  \item \textbf{Calibraci\'on}: Platt scaling opcional (\texttt{AUTOCALIBRATE\_PROBS}).
  \item \textbf{Guardado}: modelos \texttt{gat\_model\_BEST*.pt}, \texttt{GraphSMOTE\_embeddings\_model*.pt}, sidecar \texttt{\_hparams.json}.
  \item \textbf{Hash del grafo}: \texttt{\_calculate\_graph\_hash} para trazabilidad.
\end{itemize}

\section{Auditor\'ia t\'ecnica y mejores pr\'acticas}
\subsection{Fortalezas observadas}
\begin{itemize}[leftmargin=*]
  \item Separaci\'on temporal para \texttt{train/val/test} y uso de \texttt{sequence\_mask} reduce fuga temporal.
  \item GraphSMOTE en entrenamiento usa embeddings de \textbf{train} (\texttt{compute\_train\_embeddings}).
  \item \texttt{NeighborLoader} limita memoria y permite entrenamiento en grafos grandes.
  \item Mecanismos de limpieza/\emph{coalesce} para evitar inconsistencias en \texttt{edge\_attr}.
  \item Guardado de hparams y \emph{run\_id} para reproducibilidad.
\end{itemize}

\subsection{Riesgos o puntos cr\'iticos (actualizado)}
\begin{itemize}[leftmargin=*]
  \item \textbf{Fuga en balanceo manual (mitigado)}: \texttt{run\_graphsmote} ahora restringe embeddings/pool minoritario a \texttt{train\_mask}. Mantener verificaci\'on en auditor\'ias.
  \item \textbf{Dependencia de \texttt{BIDIRECCTION}}: direcci\'on de aristas sint\'eticas afecta mensajes y puede sesgar la difusi\'on de se\~nales. Se mantiene sin cambios por ahora.
  \item \textbf{Acumulaci\'on de gradientes (ajustable)}: \texttt{accumulation\_steps} ahora es configurable (config/UI). A\'un requiere tuning seg\'un memoria/hardware.
  \item \textbf{Balanceo + pesos de clase}: se deshabilita pesos con GraphSMOTE/ImGAGN (correcto), pero requiere monitoreo de m\'etricas en validaci\'on.
\end{itemize}

\subsection{Recomendaciones priorizadas (estado actual)}
\begin{enumerate}[leftmargin=*]
  \item \textbf{Implementado}: en \texttt{run\_graphsmote} (balance manual) se limita a \texttt{train\_mask} para evitar fuga potencial.
  \item \textbf{Implementado}: exponer \texttt{accumulation\_steps} en config/UI; ajustar seg\'un memoria disponible.
  \item \textbf{Pendiente/decisi\'on}: no modificar \texttt{BIDIRECCTION} por ahora; documentar el racional y evaluar impacto en AUPRC/F1 si se cambia en el futuro.
  \item \textbf{Implementado}: en ImGAGN registrar \texttt{best\_train\_recall} junto a seed y config para trazabilidad.
\end{enumerate}

\section{Anexo: referencias a funciones y archivos}
\begin{itemize}[leftmargin=*]
  \item UI principal: \texttt{src/graph\_builder\_app.py}.
  \item Construcci\'on del grafo: \texttt{src/graph.py}.
  \item Modelo GNN: \texttt{src/gat\_model.py}.
  \item Temporal head: \texttt{src/temporal\_head.py}.
  \item Entrenamiento/evaluaci\'on: \texttt{src/gnn\_main.py}, \texttt{src/train\_pretrain.py}.
  \item GraphSMOTE: \texttt{src/graphsmote.py}.
  \item ImGAGN: \texttt{src/imgagn.py}.
  \item Configuraci\'on: \texttt{src/config.py}.
\end{itemize}

\end{document}
